% Generated draft: requires mathematical and source review.
\documentclass[11pt]{article}
\usepackage[margin=25mm]{geometry}
\usepackage{fontspec}
\setmainfont{Latin Modern Roman}
\usepackage{amsmath,amssymb,mathtools,unicode-math}
\setmathfont{Latin Modern Math}
\usepackage{enumitem,needspace,xurl,hyperref}
\setlength{\parindent}{0pt}
\setlength{\parskip}{5pt}
\setlength{\emergencystretch}{3em}
\newcommand{\N}{\mathbb{N}}
\newcommand{\Z}{\mathbb{Z}}
\newcommand{\Q}{\mathbb{Q}}
\newcommand{\R}{\mathbb{R}}
\newcommand{\C}{\mathbb{C}}
\newcommand{\F}{\mathbb{F}}
\newcommand{\A}{\mathbb{A}}
\newcommand{\PP}{\mathbb{P}}
\newcommand{\E}{\mathbb{E}}
\newcommand{\eps}{\varepsilon}
\newcommand{\dd}{\,\mathrm{d}}
\begin{document}
\section*{UnsolvedMath ID: 30000523}
\textbf{Source problem code:} catalogue entry
\subsection{Definitions and mathematical statement}
For context, a finite matrix-group input consists of generators for $G\leq\operatorname{GL}_d(\F_q)$. A finite classical group is understood here as one of the linear, unitary, symplectic or orthogonal families, with the precise covering or central-quotient convention to be specified. Its natural representation is its defining action on the corresponding vector space, preserving the relevant form where applicable. Its natural or defining characteristic is the prime $\ell$ of the field defining that classical group; a representation is in natural characteristic if its ambient field also has characteristic $\ell$.

A black-box representation is an input model in which elements are encoded and accessed through specified group-operation and equality procedures, rather than through exposed matrix entries. It is an input description, and need not be an intrinsic property distinguishing an abstract group from a matrix group.

The extracted constructive-recognition task asks to identify the relevant case:
\begin{enumerate}
\item a classical group supplied in its natural representation;
\item a group supplied in another matrix representation in its natural characteristic;
\item a black-box representation;
\end{enumerate}
and to determine whether the natural characteristic is odd or even, that is, whether $\ell\ne2$ or $\ell=2$. The report's algorithmic motivation concerns standard generators, but the supplied statement does not specify required generator words or an explicit isomorphism as output.

These cases are preserved as a list of recognition settings. They are not asserted to be mutually exclusive or to exhaust all finite matrix-group inputs: a matrix group can also be accessed as a black box, and an arbitrary group need not possess a specified classical defining characteristic. A fully quantified algorithmic problem additionally requires a promise on the abstract group, an input encoding and available oracles, required output certificates, and any running-time or probabilistic correctness requirement. None is supplied here.
\subsection{Short English statement}
Determine the representation setting and the parity of the defining characteristic in constructive recognition. The source lists natural classical, other natural-characteristic matrix, and black-box settings, without giving the promises and computational model needed for a universal recognition theorem.
\subsection{Sources}
\begin{itemize}
\item Supplied report: \emph{Computational Group Theory} (2007), Oberwolfach Reports, DOI 10.4171/owr/2006/30, pp. 1795--1878; proposed year 2006. \url{https://doi.org/10.4171/owr/2006/30}.
\item Local UnsolvedMath registry, permanent ID 30000523, catalogue entry; created 25 August 2026, updated 16 September 2026. Its corrected\_verified note reports that only group notation and compound modifiers were normalized and both classification tasks retained. The supplied literature triage of 21 August 2026 labels the extraction open, describes it as a list of representation cases for standard-generator algorithms, and reports that important families are handled without a complete resolution of the extracted task. No further external reference is supplied.
\end{itemize}
\end{document}
